% =============================================================================
% Section 5 — Discussion
% =============================================================================

\section{Discussion}
\label{sec:discussion}

\subsection{Hardware Accessibility and Deployment Context}

A central design constraint of NeuroAssist AI was deployability on commodity hardware
already present in most homes. The detection pipeline requires only a~standard webcam
(720p or higher) and a~modern browser; no GPU, dedicated sensing hardware, or
proprietary dongle is needed. This stands in contrast to commercial eye-tracking AAC
devices, whose cost typically exceeds USD~3{,}000~\cite{beukelman2013aac} and whose
setup requires clinical technician involvement. A~webcam adequate for the pipeline costs
approximately BRL~100--500 (USD~20--100) and is widely available in low- and
middle-income countries (LMICs), including Brazil, where approximately 30{,}000
individuals live with MND~\cite{mehta2018mnd}. The offline-first design further
eliminates dependency on stable broadband — a~significant barrier in rural and
peri-urban settings in LMICs — since no internet connectivity is required during gesture
detection or alarm delivery on a~local network.

\subsection{Privacy-by-Design}

Facial video is classified as biometric data under both the General Data Protection
Regulation (GDPR) and Brazil's Lei Geral de Prote\c{c}\~{a}o de Dados Pessoais
(LGPD), attracting the highest level of regulatory protection for health
contexts~\cite{edpb2020privacy}. NeuroAssist AI satisfies a~strict
privacy-by-design principle: the WebAssembly runtime and the face landmarker model
execute entirely within the browser's sandboxed execution environment, and no video
frame or facial landmark coordinate is transmitted over the network. The only data that
crosses the network boundary is a~structured JSON payload containing the detected
gesture type, a~confidence score, and a~patient identifier — none of which constitutes
biometric data under current regulatory definitions. This architecture makes NeuroAssist
AI suitable for direct home deployment without requiring institutional data processing
agreements.

One honest caveat applies: the MediaPipe face landmarker model (\texttt{float16},
approximately 30\,MB) is downloaded from Google Cloud Storage on first use and
thereafter served from the browser cache. While subsequent executions require no
network access, the initial model download does involve a~third-party server. A fully
air-gapped deployment would require pre-bundling the model with the application, which
is technically straightforward but was not implemented in the current PoC.

\subsection{Extensibility via the \texttt{engineFeatureRef} Plug-in Mechanism}

The \texttt{engineFeatureRef} field in the \texttt{DetectionProfile} primitive was
designed explicitly for future gesture expansion. Each supported gesture type carries a
string token that identifies which detection algorithm the edge agent should activate
(e.g.\ \texttt{"mp.blink.double.v1"}, \texttt{"mp.brow.dual.v1"}). Adding a~new
gesture type — such as lip-corner raise, sustained eye closure, or head-nod via pose
estimation — requires only: (i) inserting a~new \texttt{TipoGesto} row with a~new
\texttt{engineFeatureRef} value, (ii) implementing the corresponding detector module in
the edge agent, and (iii) publishing a~new \texttt{DetectionProfile}. No existing
domain model, use case, or ArchUnit rule requires modification. The compile-time
enforcement of \textsc{rule-01} and \textsc{rule-02} (Table~\ref{tab:archunit})
guarantees that future contributors cannot inadvertently couple new detectors to the
domain layer.

\subsection{Limitations}

Several limitations constrain the current work and should be acknowledged explicitly.

\textbf{Illumination dependence.} EAR is computed from pixel-space landmark
coordinates; its reliability degrades in low-illumination environments where the
infrared shadow of the nasolabial fold obscures the lower eyelid landmarks. Practical
testing documented cases where the default $\theta_\text{close}=0.20$ required
manual downward adjustment to 0.17--0.18 under dim artificial lighting. Future work
should incorporate automatic illumination-adaptive threshold adjustment, for example
by normalising EAR against a~rolling baseline computed during an idle period at
session start.

\textbf{Preliminary case study scope.} The evaluation reported in this paper
includes one advanced-stage MND patient (the primary case) and five caregivers and
volunteers, and should be interpreted as a~preliminary case study rather than a~clinical
trial. The residual motor capability, blink frequency, and eyelid morphology of
advanced-stage MND patients differ substantially from those of healthy controls, and the
patient-specific calibration protocol described in Section~\ref{sec:evaluation} was
designed explicitly to accommodate this variability. Generalisation to the broader MND
population will require a~multi-site study with a~statistically representative sample
under human research ethics committee (HREC) approval.

\textbf{Polling latency overhead.} The caregiver application detects active events
by polling the \texttt{evnt-service} every five seconds. This introduces a~fixed
worst-case overhead of 5\,s to the event-to-alarm visible latency, independent of
network speed. The polling interval is a~PoC simplification; production deployment
should replace polling with server-sent events (SSE) or WebSocket push, which would
reduce this overhead to the order of tens of milliseconds.

\textbf{Stub adapters for peripheral modules.} The command catalogue (PAC service),
audit trail (HIST service), and caregiver enrolment (ACS service) are represented by
hardcoded stub adapters in the PoC. The hexagonal port interfaces are defined and
ArchUnit-enforced, but the real implementations are outside the scope of this paper.
Latency and reliability results therefore reflect only the GEST and EVNT subsystems.

\subsection{Future Work}

The most impactful next steps, in order of clinical priority, are: (i)~a~clinical
study with MND participants under HREC approval, using the protocol outlined in
Section~\ref{sec:evaluation} adapted for residual motor capability; (ii)~replacement
of the polling transport with SSE to eliminate the 5\,s alarm latency ceiling;
(iii)~integration of the full PAC command catalogue to support multi-command
communication (beyond the single emergency alert demonstrated here); and
(iv)~evaluation of additional gesture types, particularly for patients who retain
unilateral brow control after bilateral blink control is lost.

A high-priority integration pathway is coupling NeuroAssist AI distress events with NIV ventilator data streams. Research groups studying PVA in MND~\cite{stell2024pva, yang2026pva} work with multichannel physiological signals in EDF format, timestamped to second precision. NeuroAssist AI's \texttt{eventoId} carries a server-side \texttt{createdAt} timestamp and a device-side \texttt{recognizedAt} timestamp; correlating these with the NIV monitoring timeline would allow a future study to test whether caregiver-confirmed distress events co-occur with elevated PVA indices — a clinically meaningful hypothesis that neither system can test alone.